% Definitions and Technical Terminology. Built by the future pass from the
% instructions below. Do not process now. This is the bill's glossary; it must
% define every term the operative sections rely on, with USL and PSL given full
% treatment because they affect the VVUQ process synergistically.

[Write a short connecting paragraph, then render the definitions as one
\texttt{longtable} (Term, Definition, Source) so the dense glossary stays
readable across pages; do not narrate each row. Use the left-aligned ragged-right
column format for every column. Group the table into four labelled blocks. Pull
the authoritative wording from the sources in the second table below, and cite
each defining source.]

[Block 1, Statutory technical terms. Define the terms the operative sections
depend on, in plain legislative style: artificial intelligence; generative
artificial intelligence; algorithm; high-risk artificial intelligence;
AI-enabled device software function; clinical decision support; and utilization
review tool. Anchor the device-software and clinical-decision-support framings to
\texttt{cfr-devices} and \texttt{fda2023credibility}, and the utilization-review
framing to the state statutes (\texttt{ca-sb1120}, \texttt{ny-a9149},
\texttt{tx-sb1822}, \texttt{fl-hb527}).]

[Block 2, VVUQ terms. Define verification (did we build the code correctly),
validation (did we build the right code), uncertainty quantification (how
confident across repeated runs), the VVUQ gate, the verification fraction (an
equality to 1.0, not an inequality with slack), validation agreement, the
coefficient-of-variation bound, the three-way ACCEPT, BLOCK, ESCALATE decision,
and the hard catastrophe predicate. Source these verbatim in substance from
\texttt{VVUQ-02/codegen/docs/vvuq\_methodology.md} and
\texttt{vvuq\_gate\_spec.md}, and from the VVUQ-01 methods gate description. Cite
\texttt{kawchak2026vvuq01}, \texttt{kawchak2026vvuq02}, \texttt{asmevv40},
\texttt{nasastd7009a}, and \texttt{iso14971}.]

[Block 3, USL and PSL, and their synergy with VVUQ. This block is required and
must be comprehensive. Define the Unification Standard Level (USL) as the robot-
level readiness score, 1.0 to 10.0, across four equally weighted dimensions
(simulation-framework switching, AI integration, cross-robot sharing, clinical-
trial collaboration) with the five bands (Initial, Foundational, Intermediate,
Advanced, Exemplary). Define the Physical AI Standard Level (PSL) as the site-
level compliance gate across three pass/fail dimensions (legislative
authorization, regulatory compliance, operational readiness). Source both from
\texttt{psl\_usl\_standards.tex} in the national-platform final paper. Then state,
in connected prose after the table, how USL and PSL affect the VVUQ process
synergistically: PSL verifies the environment (a compliant, prepared site) and
USL validates the platform (a sufficiently mature robot), so the two form a dual-
layer verification-and-validation gate around the code-level VVUQ, and a Physical
AI trial requires all three to clear. Tie USL dimension scoring to uncertainty
quantification (dimensional variation and field-wide weakness as quantified
uncertainty) and to the Phase 0 simulation-validation pre-enrollment gate. Use the
key synergy sentence from \texttt{psl\_usl\_standards.tex} (PSL ensures the
environment is ready, USL ensures the robots are ready, a trial requires both).
Cite \texttt{kawchak2026national} and \texttt{kawchak2026sitedocs}.]

[Block 4, Physical AI system terms. Define Physical AI system, robot agent, the
robot categories and functional types and their minimum USL thresholds by
procedure type (surgical at least 7.0; therapeutic positioning and diagnostic
needle at least 6.0; rehabilitative at least 4.0; companion monitoring at least
3.0), digital twin, simulation framework, cross-framework tolerances (under 2 mm
trajectory, under 0.5 N force), emergency stop, human oversight, and the hash-
chained audit trail. Source these from
\texttt{template-paper/chunk\_02\_subpart\_a\_general\_provisions.md} (the 22
Physical AI definitions and § 312.1 thresholds) and the national MCP-PAI standard
(\texttt{kawchak2026mcp}). Keep the section symbol § for any clause reference.]

% Model sources table (parent paths at top, abbreviated leaves):
% \begin{table}[h]\centering
% \begin{tabularx}{\textwidth}{L{4.8cm} Y}
% \toprule
% \textbf{Source file (abbreviated leaf)} & \textbf{Defines} \\
% \midrule
% \multicolumn{2}{l}{\textit{Parent: papers/VVUQ-03/template-paper/}} \\
% chunk\_02\_subpart\_a\_general\_provisions.md & 22 Physical AI terms; robot categories; USL thresholds; tolerances \\
% \multicolumn{2}{l}{\textit{Parent: physical-ai-oncology-trials/national-platform/new\_paper/final\_paper/sections/}} \\
% psl\_usl\_standards.tex & USL 4 dimensions and 5 bands; PSL 3 dimensions; the synergy sentence \\
% \multicolumn{2}{l}{\textit{Parent: papers/VVUQ-02/codegen/docs/}} \\
% vvuq\_methodology.md & Verification, validation, uncertainty; model risk; context of use \\
% vvuq\_gate\_spec.md & Gate, fraction, agreement, CV bound, hard predicate \\
% \bottomrule
% \end{tabularx}
% \caption{Authoritative sources for the definitions.}
% \end{table}
